\chapter{Healthy Eating}

\section*{Definition and Overview}
Healthy eating means regularly choosing foods that provide the nutrients the body needs while limiting those that increase the risk of disease. The Australian Dietary Guidelines recommend eating a wide variety of nutritious foods from the five food groups -- vegetables, fruit, grains, lean meats and alternatives, and dairy or alternatives -- and limiting foods high in saturated fat, added sugar, and salt. Healthy eating supports growth and development, maintains a healthy weight, and reduces the risk of heart disease, type 2 diabetes, some cancers, and other chronic conditions. It is not about strict rules or deprivation but about balance, variety, and enjoying food.

\section*{Epidemiology}
Most Australians do not eat according to the dietary guidelines. Only around half of adults eat enough vegetables, and around a third eat too much discretionary (junk) food. Poor diet is a leading risk factor for chronic disease, contributing to obesity, heart disease, stroke, type 2 diabetes, and some cancers, and is responsible for a substantial proportion of preventable deaths in Australia. Improving diet quality has a large potential to improve population health.

\section*{Aetiology and Pathophysiology}
Healthy eating works through the nutrients and patterns of the diet.

\begin{itemize}
    \item \textbf{Vegetables and fruit:} provide fibre, vitamins, minerals, and protective plant compounds
    \item \textbf{Whole grains:} provide fibre and energy, and are associated with lower disease risk
    \item \textbf{Protein:} from lean meat, fish, eggs, legumes, and nuts, supports muscle and body function
    \item \textbf{Dairy or alternatives:} provide calcium and protein for bone health
    \item \textbf{Healthy fats:} unsaturated fats from nuts, seeds, avocado, and oils support heart health
    \item \textbf{Limiting harmful nutrients:} reducing saturated fat, added sugar, and salt
    \item \textbf{Energy balance:} matching energy intake to activity maintains a healthy weight
    \item \textbf{Gut health:} fibre feeds beneficial gut bacteria
\end{itemize}

\section*{Clinical Features}
Signs of a healthy diet and the effects of an unhealthy one include:

\begin{itemize}
    \item A diet with a variety of foods from all five food groups
    \item Adequate energy for activity and growth
    \item Stable, healthy weight
    \item Good energy, digestion, and wellbeing
    \item Lower risk of heart disease, diabetes, and obesity
    \item Poor diet effects: weight gain, fatigue, poor digestion, and increased disease risk
    \item Nutrient deficiencies from restricted or imbalanced diets
\end{itemize}

\section*{Differential Diagnosis}
Nutritional concerns can have medical and other causes.

\begin{itemize}
    \item Weight loss or gain from medical conditions, including thyroid disease and diabetes
    \item Nutrient deficiencies from malabsorption, including coeliac disease
    \item Food allergies and intolerances
    \item Eating disorders, which need specialist care
    \item Difficulty accessing healthy food (food insecurity)
    \item Confusing or misleading dietary information and fad diets
    \item Cultural and individual food preferences and needs
\end{itemize}

\section*{When to Seek Medical Care}
See a doctor or accredited practising dietitian for help with a medical condition affected by diet, unexplained weight change, or concerns about your eating. Most people can improve their eating with guidance from reliable sources.

\textbf{Call 000 (or local emergency services) for:}
\begin{itemize}
    \item This chapter concerns nutrition rather than emergencies; however, seek urgent care for severe dehydration, difficulty swallowing, or any serious illness
    \item Severe allergic reactions (anaphylaxis) to food
    \item Chest pain or other serious symptoms after eating in people with known conditions
\end{itemize}

\section*{Diagnosis and Investigations}
Assessment of eating involves review and, when needed, tests.

\begin{itemize}
    \item \textbf{Dietary review:} an assessment of usual food intake
    \item \textbf{Weight and body mass index:} tracking weight
    \item \textbf{Blood tests:} iron, vitamin B12, vitamin D, and other nutrients when deficiency is suspected
    \item \textbf{Medical review:} for conditions affected by diet, including diabetes and high cholesterol
    \item \textbf{Dietitian referral:} an accredited practising dietitian provides individualised advice
    \item \textbf{Follow-up:} monitoring of weight, nutrients, and health markers
\end{itemize}

\section*{Management}
Healthy eating is planned around the five food groups and individual needs.

\begin{itemize}
    \item \textbf{Eat a variety:} choose foods from all five food groups daily
    \item \textbf{Vegetables and fruit:} aim for five serves of vegetables and two of fruit daily
    \item \textbf{Whole grains:} choose wholegrain bread, cereals, rice, and pasta
    \item \textbf{Protein:} include lean meat, fish, eggs, legumes, and nuts
    \item \textbf{Dairy:} include dairy or calcium-fortified alternatives
    \item \textbf{Water:} drink water as the main drink
    \item \textbf{Limit discretionary foods:} reduce takeaway, confectionery, sugary drinks, and salty snacks
    \item \textbf{Plan meals:} planning and cooking at home supports healthy eating
    \item \textbf{Read labels:} check fat, sugar, and salt content
    \item \textbf{Individualise:} adapt guidelines to age, activity, health conditions, and preferences
\end{itemize}

\section*{Complications}
\begin{itemize}
    \item Nutrient deficiencies from restricted diets
    \item Weight gain and obesity
    \item Increased risk of heart disease, diabetes, and cancer
    \item Poor energy, mood, and sleep
    \item Dental decay from sugar
    \item Confusion and wasted spending on fad diets
    \item Eating disorders from extreme dieting
\end{itemize}

\section*{Prognosis and Outcomes}
Healthy eating has powerful effects on health. It reduces the risk of heart disease, stroke, type 2 diabetes, and some cancers, supports a healthy weight, and improves energy, mood, and longevity. Improvements in diet produce benefits at any age, and even modest changes, such as increasing vegetables and reducing sugary drinks, make a difference. With practical guidance and gradual change, most people can adopt a healthier eating pattern and maintain it for life.

\section*{Prevention and Self-Care}
\begin{itemize}
    \item Fill half your plate with vegetables and fruit
    \item Choose wholegrain over refined grains
    \item Include protein at every meal
    \item Drink water instead of sugary drinks
    \item Limit takeaway, processed foods, and snacks
    \item Cook at home and plan meals
    \item Shop with a list and read food labels
    \item Eat mindfully and without distractions
    \item Seek advice from an accredited practising dietitian for individual needs
    \item Make gradual changes that you can sustain
\end{itemize}

\section*{Sources and Further Reading}
The following reputable sources provide further information for healthcare students and patients seeking reliable, current advice about healthy eating.

\begin{itemize}
    \item Eat for Health (Australian Government). \textit{Australian Dietary Guidelines.} https://www.eatforhealth.gov.au/
    \item Dietitians Australia. \textit{Healthy eating advice.} https://dietitiansaustralia.org.au/
    \item Healthdirect Australia. \textit{Healthy eating.} https://www.healthdirect.gov.au/
    \item Australian Institute of Health and Welfare. \textit{Inadequate diet data.} https://www.aihw.gov.au/
    \item NHS. \textit{Eat well.} https://www.nhs.uk/
    \item World Health Organization. \textit{Healthy diet fact sheet.} https://www.who.int/
\end{itemize}
